% !TeX root = ../regression_logistique.tex
\chapter{Propriétés de la fonction logistique}
\label{ann:sigmoide}

\orient{%
établir $\sigma(-z)=1-\sigma(z)$ et $\sigma'=\sigma(1-\sigma)$ ; distinguer le
rôle de $\sigma'$ dans la sensibilité de la probabilité et dans le gradient}{%
règle de chaîne, dérivée d'un inverse (\annexeref{ann:prerequis})}{%
approfondissement ; le fil principal n'emploie que les deux égalités encadrées}

\section{Symétrie}

Échanger les deux réponses revient à changer le signe du score, la frontière
étant le lieu où il s'annule. La fonction logistique doit donc échanger $p$ et
$1-p$ lorsqu'on passe de $z$ à $-z$, ce que le calcul confirme.

\begin{align}
  \sigma(-z) &= \frac{1}{1+e^{-(-z)}}
    &&\text{définition, avec } -z \text{ à la place de } z \notag\\[2pt]
  &= \frac{1}{1+e^{z}}
    &&\text{car } -(-z)=z \notag\\[2pt]
\end{align}

\slidebreak

Et par ailleurs :
\begin{align}
  1-\sigma(z) &= 1-\frac{1}{1+e^{-z}}
    &&\text{définition} \notag\\[2pt]
  &= \frac{(1+e^{-z})-1}{1+e^{-z}}
    &&\text{même dénominateur} \notag\\[2pt]
  &= \frac{e^{-z}}{1+e^{-z}}
    &&\text{on simplifie le numérateur} \notag
\end{align}

\slidebreak

\begin{align}
  1-\sigma(z)
  &= \frac{e^{-z}}{1+e^{-z}}\cdot\frac{e^{z}}{e^{z}}
    &&\text{on multiplie par } 1 \text{ écrit } e^{z}/e^{z} \notag\\[2pt]
  &= \frac{e^{-z}e^{z}}{e^{z}+e^{-z}e^{z}}
    &&\text{on distribue} \notag\\[2pt]
  &= \frac{1}{e^{z}+1}
    &&\text{car } e^{-z}e^{z}=e^{0}=1 \label{eq:sym}
\end{align}

Les deux résultats coïncident, donc
\[
  \boxed{\;\sigma(-z)=1-\sigma(z)\;}
\]

La probabilité accordée à « \Gryf » pour un score $z$ égale donc celle
accordée à « pas \Gryf » pour le score $-z$.

\begin{releve}
$\sigma(2)=0{,}8808$ et $\sigma(-2)=0{,}1192$, de somme $1$.
\end{releve}

\section{Dérivée}

Écrite $\sigma=1/u$ avec $u=1+e^{-z}$, la fonction se dérive par la règle de
l'inverse, la chaîne fournissant $u'$. Le résultat s'exprime avec $\sigma$
elle-même, ce qui évite d'avoir à réévaluer une exponentielle dans le calcul du
gradient.

\begin{align}
  \sigma(z) &= \frac{1}{1+e^{-z}} = \frac{1}{u},\qquad u=1+e^{-z} \notag\\[4pt]
  u' &= \bigl(1+e^{-z}\bigr)'
     = 0+e^{-z}\cdot(-1)
    &&\text{dérivée d'une somme, puis chaîne sur } e^{-z} \notag\\[2pt]
     &= -e^{-z} \notag\\[6pt]
\end{align}

\slidebreak

\begin{align}
  \sigma'(z) &= -\frac{u'}{u^{2}}
    &&\text{dérivée d'un inverse : } (1/u)'=-u'/u^{2} \notag\\[2pt]
  &= -\frac{-e^{-z}}{(1+e^{-z})^{2}}
    &&\text{on remplace } u \text{ et } u' \notag\\[2pt]
  &= \frac{e^{-z}}{(1+e^{-z})^{2}}
    &&\text{les deux signes moins s'annulent} \notag\\[2pt]
  &= \frac{1}{1+e^{-z}}\cdot\frac{e^{-z}}{1+e^{-z}}
    &&\text{on sépare le carré en deux facteurs} \notag\\[2pt]
  &= \sigma(z)\cdot\bigl(1-\sigma(z)\bigr)
    &&\text{le second facteur est } 1-\sigma(z),\ \text{établi en \eqref{eq:sym}}
\end{align}

\begin{resultat}[Dérivée de la sigmoïde]
\[
  \sigma'(z)=\sigma(z)\bigl(1-\sigma(z)\bigr).
\]
La pente se calcule donc sans nouvelle exponentielle : si l'on connaît déjà
$p=\sigma(z)$, la pente vaut $p(1-p)$.
\end{resultat}

\section{Variation de la probabilité avec le score}

La quantité $p(1-p)$ atteint son maximum, $0{,}25$, en $p=0{,}5$, et décroît vers
$0$ aux deux extrémités : la probabilité varie fortement avec le score près de la
frontière, et très peu lorsqu'elle est déjà proche de $0$ ou de $1$.

\begin{figure}[htbp]
\centering
\begin{graphique}{Deux courbes superposées : la fonction logistique croissante en S, et sa dérivée en cloche symétrique, maximale à 0,25 à l'origine et tendant vers zéro des deux côtés.}
\begin{tikzpicture}
\begin{axis}[courseplot,height=6cm,axis lines=middle,
  xlabel={score $z$},ylabel={},xmin=-6.5,xmax=6.5,ymin=-.08,ymax=1.1,
  samples=200,domain=-6.5:6.5,
  legend style={at={(.5,1.03)},anchor=south,legend columns=2,draw=none}]
  \addplot[accent,very thick]{1/(1+exp(-x))};
  \addlegendentry{$\sigma(z)$}
  \addplot[warning,very thick]{exp(-x)/(1+exp(-x))^2};
  \addlegendentry{$\sigma'(z)=\sigma(z)(1-\sigma(z))$}
  \addplot[only marks,mark=*,ink,mark size=2.4pt] coordinates {(0,0.25)};
  \node[ink,anchor=south west,font=\small] at (axis cs:.2,.25){$0{,}25$};
\end{axis}
\end{tikzpicture}
\end{graphique}
\caption{La dérivée atteint son maximum, $0{,}25$, en $z=0$, et tend vers zéro dès
que la probabilité approche $0$ ou $1$.}
\label{fig:derivee-sigmoide}
\end{figure}
\FloatBarrier

Cette cloche gouverne deux quantités. La première est la sensibilité de la
\emph{probabilité} au score : $p$ varie fortement avec $z$ près de la frontière
et très peu une fois proche de $0$ ou de $1$, ce qui explique qu'agrandir la
norme des coefficients modifie les probabilités sans déplacer la frontière
(\annexeref{ann:norme}). La seconde est la \emph{courbure} de la perte, dont
$\sigma'$ fournit le facteur $s_i$ de la hessienne
(\annexeref{ann:convexite}).

Le poids d'une observation dans le gradient est en revanche
$\partial\ell/\partial z=p-y$ (\annexeref{ann:gradient}) : le facteur
$\sigma(1-\sigma)$ apparaît dans la dérivation, puis se simplifie avec les
dénominateurs venus du logarithme. Les deux quantités se séparent nettement. Sur
la frontière, $p=0{,}5$ : $\sigma'$ atteint son maximum $\frac14$ tandis que
$|p-y|$ vaut $0{,}5$ ; pour une observation d'étiquette $y=1$ classée à
$p=0{,}02$, $\sigma'$ tombe à $0{,}0196$ tandis que $|p-y|$ monte à $0{,}98$. La
correction est donc portée par les prédictions incorrectes et assurées.
